\section{How much data does inference require?}\label{sec:cost}
Even with perfect detector knowledge, dependence can leave a short record uninformative. Correlated samples can repeat nearly the same information, so an informative frequency band may require a long record. We first remove calibration and source-fitting uncertainty to isolate the recording cost imposed by the source itself.

Consider unit-variance Gaussian channels with identity responses, so the sensors leave the source unchanged. Their marginal spectral densities lie between $1/2$ and $\mathcal S/2$. The spectral dynamic range $\mathcal S\ge3$ measures the ceiling relative to the floor. We seek an absolute covariance contrast at least $\gamma$, with $0<\gamma\le1/4$. A difficult alternative places its directional signal in two narrow frequency bands. A length-$N$ record resolves only about $N/(\mathcal S-1)$ directions in those bands. Many samples then carry little new evidence about the asymmetry. Relative entropy measures how distinguishable this record's probability law is from a reversible law and therefore limits every test. The same information measure relates forward and reversed trajectories to dissipation in stochastic thermodynamics~\cite{roldan2010,seifert2012}.

\begin{theorem}[Recording cost]\label{thm:main-cost}
For each $(\mathcal S,\gamma)$ in this class, SM Sec.~S6 constructs a fixed reversible source and a fixed alternative with contrast $-\gamma$ and the same known marginal spectra. The necessary length inequality in SM Sec.~S6.3 applies to any test of this pair with false-rejection probability $\alpha$ and detection probability $1-\beta$, where $0<\alpha,\beta<1$ and $\alpha+\beta<1$. For the full class, the calibrated test with exact identity responses and $q=\mathcal S$ has false-rejection probability at most $0.05$ and detection probability at least $0.95$ whenever
\begin{equation}
 N\ge1024\mathcal S/\gamma^2.\label{eq:cost-uniform}
\end{equation}
At $\mathcal S=9$, $\gamma=0.1$, and $\alpha=\beta=0.05$, at least $5214$ samples are necessary for the fixed pair; direct evaluation of the sharper sufficient condition gives $129580$ samples for the full class.
\end{theorem}
Thus dependence can force thousands of samples even with exact responses. SM Sec.~S6, Theorems~S8--S9, gives the complete necessary inequality, its finite trace-defect bound, and the sufficient calculation. The information required by the two error rates is $(1-\beta)\log[(1-\beta)/\alpha]+\beta\log[\beta/(1-\alpha)]$. A record is necessarily too short if an upper bound on its available relative entropy is smaller. The sufficient-to-necessary length ratio is below $25$ in the worked case. It compares a class-wide sufficient bound with a fixed-pair lower bound, so it is not an optimal-length claim.
